
\section{\BookRecordIndex{T0047}Seidel 表示}\label{sec:11-4}

本节为满足强\BookRecordIndex{T0058}半正性条件~\eqref{eq:8-5-1} 的\BookRecordIndex{T0358}辛流形定义 \BookRecordIndex{T0047}Seidel 表示。我们说明如何用它
理解一般\BookRecordIndex{T0169}环面流形的\BookRecordIndex{T0229}量子上同调环结构。作为另一个应用，我们还将计算任意\BookRecordIndex{T0156}哈密顿纤维化
$\widetilde M\to S^2$ 总空间的加性上同调，其中纤维是\BookRecordIndex{T0058}半正\BookRecordIndex{T0358}辛流形 $(M,\omega)$
（命题~\ref{prop:11-4-5}）。

在例~\ref{ex:8-6-8} 中，我们讨论了单调情形下的 \BookRecordIndex{T0047}Seidel 表示。在该情形中，可以把哈密顿
环路 $\psi$ 的不变量 $\mathcal S(\psi)$ 定义为 $M$ 的一个上同调元素。式~\eqref{eq:8-6-4}
中不变量 $\mathcal S_{\widetilde A}(\psi)$ 的定义原封不动地适用于\BookRecordIndex{T0058}半正情形。不过，现在可能有
无穷多个\BookRecordIndex{T0189}截面类 $\widetilde A$ 具有非零不变量 $\mathcal S_{\widetilde A}(\psi)$，而
$\mathcal S(\psi)$ 将定义为\BookRecordIndex{T0229}量子上同调中的一个（可逆）元素。下面给出细节。

设 $(M,\omega)$ 是满足~\eqref{eq:8-5-1} 的闭\BookRecordIndex{T0358}辛流形。则与 $M$ 上一般二参数\BookRecordIndex{T0088}殆复结构族
相联系的每条 \BookRecordIndex{T0003}$J$-全纯曲线都有非负陈数。令
$\psi=\{\psi_t\}_{t\in\mathbb R/\mathbb Z}$ 是 $M$ 的\BookRecordIndex{T0161}哈密顿辛微分同胚环路，并令
$\widetilde M_\psi\to S^2$ 是式~\eqref{eq:8-2-6} 中以 $M$ 的辛同胚副本为纤维的\BookRecordIndex{T0156}哈密顿
纤维化。由推论~\ref{cor:9-1-2}，对每个 $x\in M$，环路
$t\mapsto\psi_t(x)$ 都可缩。这蕴含\BookRecordIndex{T0279}纤维化 $\widetilde M_\psi$ 有截面。给定可由截面表示的
同调类 $\widetilde A\in H_2(\widetilde M_\psi)$，如式~\eqref{eq:8-6-4} 那样定义上同调类
\begin{equation}\label{eq:11-4-1}
  \mathcal S_{\widetilde A}(\psi)
  :=\sum_{\nu,\mu}
  \operatorname{GW}^{\widetilde M_\psi,\mathbf w}_{\widetilde A,1}(e_\nu)
  g^{\nu\mu}e_\mu
  \in H^{-2c}(M).
\end{equation}
这里 $c:=c_1^{\mathrm{Vert}}(\widetilde A)\le0$ 是 $\widetilde A$ 的竖直陈数，
$e_0,\ldots,e_N$\sourcefootnote{原书把全上同调基的末指标印作 $n$；但同句中的
$n=\dim_{\mathbb R}M/2$，而本章固定的上同调基是 $e_0,\ldots,e_N$，故补正为 $N$。}
是 $H^*(M)$ 的一个基，而 $g^{\nu\mu}$ 是矩阵
\[
  g_{\nu\mu}:=\int_M e_\nu\smile e_\mu
\]
的逆矩阵。\BookRecordIndex{T0143}格罗莫夫--威滕不变量
\[
  \operatorname{GW}^{\widetilde M_\psi,\mathbf w}_{\widetilde A,1}:
  H^{2n+2c}(M)\longrightarrow\mathbb Z
\]
按定义~\ref{def:8-6-6} 理解，其中有一个\BookRecordIndex{T0150}固定标记点 $w\in S^2$。具体地，通过嵌入
$\iota:M\to\widetilde M_\psi$，把 $M$ 与 $\widetilde M_\psi$ 在 $w\in S^2$ 上的纤维
等同；于是 $\mathcal S_{\widetilde A}(\psi)$ 是\BookRecordIndex{T0337}伪循环
\[
  \iota^{-1}\circ\widetilde{\operatorname{ev}}^{w,J,H}:
  \mathcal M(\widetilde A;J,H)\longrightarrow M
\]
的 Poincar\'e 对偶上同调类。该\BookRecordIndex{T0337}伪循环定义在 $\widetilde M_\psi$ 中类 $\widetilde A$ 的
$J$-全纯截面\BookRecordIndex{T0246}模空间上。从几何上说，它在子流形 $X\subset M$ 上的积分，是穿过
$\iota(X)$ 的类 $\widetilde A$ 的 $J$-全纯截面数，并计入适当的符号。因为嵌入 $\iota$
在 $M$ 的\BookRecordIndex{T0158}哈密顿同痕意义下唯一，上同调类 $\mathcal S_{\widetilde A}(\psi)$ 与 $\iota$ 的
选择无关。下面由这些上同调类构造\BookRecordIndex{T0229}量子上同调类。

\begin{definition}\label{def:11-4-1}
设 $(M,\omega)$ 是满足~\eqref{eq:8-5-1} 的闭\BookRecordIndex{T0358}辛流形，
$\psi=\{\psi_t\}_{t\in\mathbb R/\mathbb Z}$ 是\BookRecordIndex{T0161}哈密顿辛微分同胚环路，而
$\sigma:S^2\to\widetilde M_\psi$ 是一个截面。固定定义~\ref{def:11-1-3} 中任意一个\BookRecordIndex{T0231}量子
系数环 $\Lambda$。定义 \emph{\BookRecordIndex{T0046}Seidel 元}
\[
  \mathcal S(\psi,\sigma)\in\operatorname{QH}^*(M;\Lambda)
\]
为
\begin{equation}\label{eq:11-4-2}
  \mathcal S(\psi,\sigma)
  :=\sum_{A\in H_2(M)}
  \mathcal S_{[\sigma]+\iota_*A}(\psi)\otimes e^A.
\end{equation}
\end{definition}

作为例子，考虑环 $\Lambda:=\Lambda^{\mathrm{univ}}[q,q^{-1}]$，即以\BookRecordIndex{T0270}普适 Novikov 环为
系数的 $q,q^{-1}$ 多项式环（见例~\ref{ex:11-1-4}(v)）。变量 $q$ 的次数为二。借助
Guillemin--Lerman--Sternberg \srcite{163} 的耦合类
$u_\psi\in H^2(\widetilde M_\psi)$，便可消去 \BookRecordIndex{T0046}Seidel 元对截面 $\sigma$ 的依赖。该类定义为
闭二形式 $\widetilde\omega\in\Omega^2(\widetilde M_\psi)$ 的上同调类，其中
$\widetilde\omega$ 在每个纤维上的限制都是给定\BookRecordIndex{T0363}辛形式，并且
$\widetilde\omega^{n+1}$ 的\BookRecordIndex{T0374}沿纤维积分为零。其上同调类与 $\widetilde\omega$ 的选择无关
（见命题~\ref{prop:8-2-2}）。

\begin{definition}\label{def:11-4-2}
设 $(M,\omega)$ 和 $\psi$ 如定义~\ref{def:11-4-1}，并记
\[
  \Lambda:=\Lambda^{\mathrm{univ}}[q,q^{-1}].
\]
定义 \BookRecordIndex{T0046}Seidel 元 $\mathcal S(\psi)\in\operatorname{QH}^*(M;\Lambda)$ 为
\begin{equation}\label{eq:11-4-3}
  \mathcal S(\psi)
  :=\sum_{\widetilde A}
  \mathcal S_{\widetilde A}(\psi)
  \otimes q^{c(\widetilde A)}t^{\widetilde\omega(\widetilde A)}.
\end{equation}
这里 $\widetilde\omega\in\Omega^2(\widetilde M_\psi)$ 是 $\widetilde M_\psi$ 上一个
\BookRecordIndex{T0155}哈密顿联络的\BookRecordIndex{T0260}耦合形式，而 $c:=c_1^{\mathrm{Vert}}\in H^2(\widetilde M_\psi)$ 是竖直第一
陈类。求和遍历所有可由 $\widetilde M_\psi$ 的截面表示、并满足
$-\dim M\le2c(\widetilde A)\le0$ 的同调类
$\widetilde A\in H_2(\widetilde M_\psi)$。
\end{definition}

定义~\ref{def:11-4-2} 中 \BookRecordIndex{T0047}Seidel 表示的表述，是 Lalonde--McDuff--Polterovich
\srcite{228} 所用表述的一个变体：他们把 $\Lambda^{\mathrm{univ}}[q,q^{-1}]$ 换成
Novikov 环 $\Lambda_\omega$，其底群是球面同调类群 $H_2^S(M)$ 关于下述等价关系的商：
$A\sim B$ 当且仅当 $\omega(A)=\omega(B)$ 且 $c_1(A)=c_1(B)$。在嵌入
\[
  e^A\longmapsto q^{c_1(A)}t^{\omega(A)}
\]
下，这个 Novikov 环是 $\Lambda^{\mathrm{univ}}[q,q^{-1}]$ 的子环。换言之，把同调群
$H_2(M)$ 除以同态 $A\mapsto(\omega(A),c_1(A))$ 的核，等价于只保留同调类 $A$ 的
$\omega(A)$ 与 $c_1(A)$ 信息。

\begin{theorem}[Seidel]\label{thm:11-4-3}
设 $(M,\omega)$ 是满足\BookRecordIndex{T0058}半正性条件~\eqref{eq:8-5-1} 的闭\BookRecordIndex{T0358}辛流形，并令
$\Lambda:=\Lambda^{\mathrm{univ}}[q,q^{-1}]$。则对所有
$\phi,\psi\in\pi_1(\operatorname{Ham}(M,\omega))$，都有
\[
  \mathcal S(\phi\psi)=\mathcal S(\phi)*\mathcal S(\psi),
  \qquad
  \mathcal S(\operatorname{id})=1.
\]
因此，对每个 $\psi\in\pi_1(\operatorname{Ham}(M,\omega))$，\BookRecordIndex{T0046}Seidel 元
$\mathcal S(\psi)\in\operatorname{QH}^*(M;\Lambda)$ 都属于
$\operatorname{QH}^*(M;\Lambda)$ 的偶次乘法单位群，而且映射
$\psi\mapsto\mathcal S(\psi)$ 是群同态。
\end{theorem}

断言 $\mathcal S(\operatorname{id})=1$ 是关于乘积流形
$\widetilde M:=S^2\times M$ 的\BookRecordIndex{T0143}格罗莫夫--威滕不变量的陈述，由习题~\ref{ex:7-3-4}
成立。恒等式
$\mathcal S(\phi\psi)=\mathcal S(\phi)*\mathcal S(\psi)$ 可重述为\BookRecordIndex{T0156}哈密顿纤维化的
\BookRecordIndex{T0143}格罗莫夫--威滕不变量所满足的（分裂）公理的类似物；其证明仍基于第 10 章的\BookRecordIndex{T0254}粘合论证。
虽然这里给出的证明只适用于\BookRecordIndex{T0058}半正情形，该定理事实上对所有闭\BookRecordIndex{T0358}辛流形均成立；见 McDuff
\srcite{269}。还要注意，在量子乘法等于通常杯积的流形上，例如当 $[\omega]$ 在
$\pi_2(M)$ 上为零时，上述恒等式退化为定理~\ref{thm:8-6-9} 中的断言
$\mathcal S(\phi\psi)=\mathcal S(\phi)\smile\mathcal S(\psi)$。

\begin{remark}\label{rem:11-4-4}
\BookRecordIndex{T0046}Seidel 元 $\mathcal S(\phi)$ 通过
\[
  \operatorname{QH}^*(M;\Lambda)\longrightarrow
  \operatorname{QH}^*(M;\Lambda),
  \qquad a\longmapsto\mathcal S(\phi)*a
\]
诱导\BookRecordIndex{T0229}量子上同调的一个自同构，也记作 $\mathcal S(\phi)$。我们将在下面
定理~\ref{thm:11-4-3} 的证明中看到
\begin{equation}\label{eq:11-4-4}
  \mathcal S(\phi)*a
  =\sum_{\widetilde A,\nu,\mu}
  \operatorname{GW}^{\widetilde M_\phi,\mathbf w}_{\widetilde A,2}(a,e_\nu)
  g^{\nu\mu}e_\mu\otimes
  q^{c(\widetilde A)}t^{\widetilde\omega(\widetilde A)}.
\end{equation}
所以它由带两个\BookRecordIndex{T0150}固定标记点 $w_1,w_2$ 的截面不变量给出。这里应把类 $a$ 的 Poincar\'e
对偶看成由 $w_1\in S^2$ 上纤维 $M_1$ 中的循环表示，而把 $e_\nu$ 的 Poincar\'e 对偶
看成由 $w_2\in S^2$ 上纤维 $M_2$ 中的循环表示。不变量
$\operatorname{GW}^{\widetilde M,\mathbf w}_{\widetilde A,2}(a,e_\nu)$ 计数穿过这些循环的
$\widetilde M$ 中 $J$-全纯截面。由于每个纤维与 $M$ 的等同在\BookRecordIndex{T0161}哈密顿辛微分同胚意义下
良定义，$\mathcal S(\phi)$ 的确给出 $\operatorname{QH}^*(M)$ 的良定义自同构。

用同调语言说，给定由纤维 $M_1$ 中一个循环（也记作 $\alpha$）表示的类
$\alpha\in H_*(M)$，类 $\mathcal S(\phi)(\alpha)$ 是诸项
$\beta_{\widetilde A}\otimes q^{-c(\widetilde A)}t^{-\widetilde\omega(\widetilde A)}$ 之和，
其中同调类 $\beta_{\widetilde A}$ 由所有穿过 $\alpha$ 的类 $\widetilde A$ 截面与 $M_2$
的交表示。在这一解释下，\BookRecordIndex{T0046}Seidel 元
$\mathcal S(\phi):=\mathcal S(\phi)(1)$ 是单位元的像，而
定理~\ref{thm:11-4-3} 断言 $\mathcal S(\phi\psi)$ 是这些自同构的复合
$\mathcal S(\phi)\circ\mathcal S(\psi)$。因此群
$\pi_1(\operatorname{Ham}(M,\omega))$ 通过左右
$\operatorname{QH}^*(M)$-模同构作用在 $M$ 的\BookRecordIndex{T0229}量子上同调上，即
\[
  \mathcal S(\phi)*(a*b)
  =a*\bigl(\mathcal S(\phi)(b)\bigr)
  =\mathcal S(\phi)(a)*b,
  \qquad a,b\in\operatorname{QH}^*(M;\Lambda).
\]
\sourcefootnote{原书末项印作未定义的 $\mathcal S(a)*b$；本段定义的算子是
$\mathcal S(\phi)$，且左右模同态性质唯一要求此处为 $\mathcal S(\phi)(a)*b$。}
最近 Hutchings \srcite{194} 发现了这一表示的推广：它定义一系列群同态，从
$\pi_{2k+1}(\operatorname{Ham}(M,\omega))$ 到 $\operatorname{QH}^*(M)$ 的
$2k$ 次 $\operatorname{QH}^*(M)$-模自同态所成的加法群
$\operatorname{End}_{2k}(\operatorname{QH}^*)$。
\end{remark}

\BookRecordIndex{T0047}Seidel 表示有许多有趣应用。我们已经见过两个：在第~\ref{sec:8-6} 节中，我们用它考察
\BookRecordIndex{T0161}哈密顿辛微分同胚环路的辛作用量（见推论~\ref{cor:8-6-10}）；命题~\ref{prop:9-6-4}
则用它估计环路的 Hofer 长度。它也是 McDuff 关于单有理流形的工作 \srcite{273} 中的
关键成分，因为它能保证某些\BookRecordIndex{T0143}格罗莫夫--威滕不变量不消失。正如第 12 章将看到的，它是
哈密顿群上\BookRecordIndex{T0132}弗洛尔同调丛的\BookRecordIndex{T0095}单值化，因而也是理解\BookRecordIndex{T0271}谱不变量的重要工具（参见
\srcite{307,105,274}）。它对\BookRecordIndex{T0156}哈密顿纤维化也有推论。这里只给出取自
Lalonde--McDuff--Polterovich \srcite{228} 的最简单推论；这个方向上的进一步结果见
\srcite{228} 及 Lalonde--McDuff \srcite{227}。

\begin{proposition}\label{prop:11-4-5}
设 $(M,\omega)$ 是满足条件~\eqref{eq:8-5-1} 的闭\BookRecordIndex{T0358}辛流形，并令
$\pi:\widetilde M\to S^2$ 是纤维辛同胚于 $M$ 的\BookRecordIndex{T0156}哈密顿纤维化。则映射
\sourcefootnote{原书显示的映射不可能是所称同构：按紧随其后的定义，$\iota_!$ 把次数提高
二，而且投影公式给出 $\pi^*[S^2]\smile\iota_!b=0$。本证明实际由 $\iota_*$ 单射\BookRecordIndex{T0326}推出
$\iota^*$ 满射，从而在任选一个线性提升 $s:H^*(M)\to H^*(\widetilde M)$ 后得到
$a\otimes b\mapsto\pi^*a\smile s(b)$ 的非典范加性同构；原书未指定这样的提升，故此处
保留原显示式。}
\[
  H^*(S^2)\otimes H^*(M)\longrightarrow H^*(\widetilde M),
  \qquad a\otimes b\longmapsto\pi^*a\smile\iota_!b
\]
在有理上同调上是同构。这里 $\iota:M\to\widetilde M$ 表示一个纤维的包含映射，而
$\iota_!$ 是由通常的同调推前 $\iota_*$ 定义的上同调推前：
\[
  \iota_!:=\operatorname{PD}_{\widetilde M}\circ\iota_*\circ\operatorname{PD}_M.
\]
\end{proposition}

\begin{proof}
每个\BookRecordIndex{T0156}哈密顿纤维化 $\widetilde M\to S^2$ 都像式~\eqref{eq:8-2-6} 那样，由一个哈密顿环路
$\phi$ 把两份 $\mathbb C\times M$ 粘在一起而定义。因此，可把
$\widetilde M\to S^2$ 等同为形如 $\widetilde M_\phi\to S^2$ 的\BookRecordIndex{T0279}纤维化。只须证明纤维的
包含映射 $\iota:M\to\widetilde M_\phi$ 在有理同调上诱导单射。该结论立即由
定理~\ref{thm:11-4-3} 得到。因为映射
$a\mapsto\mathcal S(\phi)*a$ 是\BookRecordIndex{T0229}量子上同调的自同构，所以
$\mathcal S(\phi)*a\ne0$ 当且仅当 $a\ne0$。由式~\eqref{eq:11-4-4}，
$\mathcal S(\phi)*a$ 是一个和，其系数是形如
$\operatorname{GW}_{\widetilde A,2}(a,e_\nu)$ 的\BookRecordIndex{T0143}格罗莫夫--威滕不变量；除非对偶同调类
$\operatorname{PD}(a)\in H_*(M)$ 在 $H_*(\widetilde M)$ 中的像
$\iota_*\operatorname{PD}(a)$ 非零，这些不变量都消失。这就证明了命题。
\end{proof}

现在转向定理~\ref{thm:11-4-3} 的证明。它以 $2$-球面上\BookRecordIndex{T0156}哈密顿纤维化的
\BookRecordIndex{T0143}格罗莫夫--威滕不变量所满足的（分裂）公理类似物为基础。为表述该结果，先观察：给定两个
哈密顿环路 $\phi,\psi\in\pi_1(\operatorname{Ham}(M,\omega))$，由逐点乘积
$\{\phi_t\psi_t\}$ 确定的\BookRecordIndex{T0279}纤维化 $\widetilde M_{\phi\psi}$ 可与\BookRecordIndex{T0278}纤维和
$\widetilde M_\phi\#\widetilde M_\psi$ 等同。它提升了底球面上通常的\BookRecordIndex{T0227}连通和运算
$S^2\#S^2=S^2$：从 $\widetilde M_\phi,\widetilde M_\psi$ 各自切去一个纤维 $F$ 的邻域，
再沿边界适当地\BookRecordIndex{T0254}粘合余下两块。（准确公式见下文。还要注意，$\phi,\psi$ 的次序无关紧要，
因为环路 $\{\phi_t\psi_t\}$ 与 $\{\psi_t\phi_t\}$ 同伦；所以
$\widetilde M_{\phi\psi}=\widetilde M_{\psi\phi}$。）相应地，在 $\widetilde M$ 中可由截面
表示的同调类集合 $H_2^\sigma(\widetilde M)$ 上有\BookRecordIndex{T0227}连通和运算：取穿过被删纤维上对应点的
截面来表示这些类，再取其\BookRecordIndex{T0227}连通和所表示的类。把该映射记作
\[
  H_2^\sigma(\widetilde M_\phi)\times
  H_2^\sigma(\widetilde M_\psi)
  \longrightarrow H_2^\sigma(\widetilde M_{\phi\psi}),
  \qquad
  (\widetilde A,\widetilde B)\longmapsto\widetilde A\#\widetilde B.
\]
它满足
\begin{equation}\label{eq:11-4-5}
  c(\widetilde A\#\widetilde B)=c(\widetilde A)+c(\widetilde B),
  \qquad
  \widetilde\omega(\widetilde A\#\widetilde B)
  =\widetilde\omega(\widetilde A)+\widetilde\omega(\widetilde B).
\end{equation}

\begin{theorem}\label{thm:11-4-6}
设 $(M,\omega)$ 是满足~\eqref{eq:8-5-1} 的闭\BookRecordIndex{T0358}辛流形，并令
$\phi,\psi\in\pi_1(\operatorname{Ham}(M,\omega))$。则对所有
$\widetilde C\in H_2^\sigma(\widetilde M_{\phi\psi})$、所有
$a_1,\ldots,a_k\in H^*(M)$ 和任意整数 $0\le\ell\le k$，都有
\[
\begin{aligned}
  \operatorname{GW}^{\widetilde M_{\phi\psi}}_{\widetilde C,k}
  (a_1,\ldots,a_k)
  ={}&\sum_{\widetilde A\#\widetilde B=\widetilde C}
  \sum_{\nu,\mu}
  \operatorname{GW}^{\widetilde M_\phi}_{\widetilde A,\ell+1}
  (a_1,\ldots,a_\ell,e_\nu)g^{\nu\mu}\\
  &\qquad\qquad\cdot
  \operatorname{GW}^{\widetilde M_\psi}_{\widetilde B,k-\ell+1}
  (e_\mu,a_{\ell+1},\ldots,a_k).
\end{aligned}
\]
这里所有不变量均按具有\BookRecordIndex{T0150}固定标记点来理解，即
\[
  \operatorname{GW}^{\widetilde M}_{\widetilde A,j}
  :=\operatorname{GW}^{\widetilde M,\mathbf w}_{\widetilde A,j},
  \qquad \mathbf w:=\{w_1,\ldots,w_j\}.
\]
\end{theorem}

这是描述任意\BookRecordIndex{T0278}纤维和之\BookRecordIndex{T0143}格罗莫夫--威滕不变量的一般公式的一个很简单的特例；辛情形的一般
公式由 Ionel--Parker \srcite{201} 给出。

\begin{proof}[定理~\ref{thm:11-4-3} 的证明]
在定理~\ref{thm:11-4-6} 中取 $\psi_t\equiv\operatorname{id}$ 和 $k=2$，便有
\begin{equation}\label{eq:11-4-6}
  \operatorname{GW}^{\widetilde M_\phi}_{\widetilde A,2}(a,b)
  =\sum_{\widetilde B+\iota_*C=\widetilde A}
   \sum_{i,j}
   \operatorname{GW}^{\widetilde M_\phi}_{\widetilde B,1}(e_i)
   g^{ij}\operatorname{GW}^M_{C,3}(e_j,a,b)
\end{equation}
其中 $\widetilde A\in H_2^S(\widetilde M_\phi)$
\sourcefootnote{原书把这里的总空间印作 $\widetilde M_\psi$；但本步已取
$\psi_t\equiv\operatorname{id}$，式~\eqref{eq:11-4-6} 的左端、求和条件及类关系全都位于
$\widetilde M_\phi$，故补正下标为 $\phi$。}，且 $a,b\in H^*(M)$。这里用
式~\eqref{eq:8-6-2} 和定义~\ref{def:7-3-7} 把不变量
$\operatorname{GW}^{\widetilde M_\psi}_{\widetilde C,k}$ 等同于
$\operatorname{GW}^M_{C,k}$，又用式~\eqref{eq:11-4-5} 把 $\widetilde A$ 等同于
$\widetilde B+\iota_*C$。从现在起，$\operatorname{GW}$ 上标中的 $\mathbf w$ 均省略不写。
于是
\[
\begin{aligned}
  \mathcal S(\phi)*a
  &={}
  \sum_{\widetilde B,C,i,j}
  \operatorname{GW}^{\widetilde M_\phi}_{\widetilde B,1}(e_i)g^{ij}
  (e_j*a)_C\otimes
  q^{c(\widetilde B+\iota_*C)}
  t^{\widetilde\omega(\widetilde B+\iota_*C)}\\
  &={}
  \sum_{\widetilde B,C,i,j,\nu,\mu}
  \operatorname{GW}^{\widetilde M_\phi}_{\widetilde B,1}(e_i)g^{ij}
  \operatorname{GW}^M_{C,3}(e_j,a,e_\nu)g^{\nu\mu}e_\mu\otimes
  q^{c(\widetilde B+\iota_*C)}
  t^{\widetilde\omega(\widetilde B+\iota_*C)}\\
  &={}
  \sum_{\widetilde A,\nu,\mu}
  \operatorname{GW}^{\widetilde M_\phi}_{\widetilde A,2}(a,e_\nu)
  g^{\nu\mu}e_\mu\otimes
  q^{c(\widetilde A)}t^{\widetilde\omega(\widetilde A)}.
\end{aligned}
\]
最后一个等号由式~\eqref{eq:11-4-6} 得到。因此我们证明了式~\eqref{eq:11-4-4}。
进而
\[
\begin{aligned}
  \mathcal S(\phi)*\mathcal S(\psi)
  &={}
  \sum_{\widetilde A,\widetilde B,\nu,\mu}
  \operatorname{GW}^{\widetilde M_\phi}_{\widetilde A,2}
  \bigl(\mathcal S_{\widetilde B}(\psi),e_\nu\bigr)
  g^{\nu\mu}e_\mu\otimes
  q^{c(\widetilde A)+c(\widetilde B)}
  t^{\widetilde\omega(\widetilde A)+\widetilde\omega(\widetilde B)}\\
  &={}
  \sum_{\widetilde A,\widetilde B,\nu,\mu,i,j}
  \operatorname{GW}^{\widetilde M_\psi}_{\widetilde B,1}(e_i)g^{ij}
  \operatorname{GW}^{\widetilde M_\phi}_{\widetilde A,2}(e_j,e_\nu)g^{\nu\mu}e_\mu
  \otimes q^{c(\widetilde A\#\widetilde B)}
  t^{\widetilde\omega(\widetilde A\#\widetilde B)}\\
  &={}
  \sum_{\widetilde C,\nu,\mu}
  \operatorname{GW}^{\widetilde M_{\phi\psi}}_{\widetilde C,1}(e_\nu)
  g^{\nu\mu}e_\mu\otimes
  q^{c(\widetilde C)}t^{\widetilde\omega(\widetilde C)}\\
  &=\mathcal S(\phi\psi).
\end{aligned}
\]
这里第二个等号使用了式~\eqref{eq:11-4-5} 和 $\mathcal S_{\widetilde B}(\psi)$ 的公式，
第三个等号则来自取 $k=1$ 的定理~\ref{thm:11-4-6}。由于
$\psi\phi=\phi\psi\in\pi_1(\operatorname{Ham}(M))$，各项的次序无关紧要。这就证明了
定理~\ref{thm:11-4-3}。
\end{proof}

\begin{proof}[定理~\ref{thm:11-4-6} 的证明]
固定 $M$ 的两个\BookRecordIndex{T0161}哈密顿辛微分同胚环路
\[
  \phi^0=\{\phi_t^0\}_{t\in\mathbb R/\mathbb Z},
  \qquad
  \phi^\infty=\{\phi_t^\infty\}_{t\in\mathbb R/\mathbb Z}.
\]
证明的思路，是把第 10 章的\BookRecordIndex{T0256}粘合定理推广到\BookRecordIndex{T0156}哈密顿纤维化。若\BookRecordIndex{T0279}纤维化
\[
  \widetilde M^0:=\widetilde M_{\phi^0},
  \qquad
  \widetilde M^\infty:=\widetilde M_{\phi^\infty}
\]
的两个 $J$-全纯截面在一对纤维中相交（这两个纤维都与 $M$ 等同），并满足适当的\BookRecordIndex{T0166}横截性
条件，就可以把它们\BookRecordIndex{T0254}粘合成\BookRecordIndex{T0278}纤维和
\[
  \widetilde M^{\infty0}:=\widetilde M_{\phi^\infty\phi^0}
\]
的一个 $J$-全纯截面。理解设定和\BookRecordIndex{T0394}预粘合构造之后，第 10 章的证明便逐字适用于当前情形，
所以我们只描述主要思路。

根据式~\eqref{eq:8-2-6} 中的定义，\BookRecordIndex{T0156}哈密顿纤维化
$\widetilde M^0\to S^2$ 是商空间
\[
  \widetilde M^0=
  \frac{\{\pm1\}\times\mathbb C\times M}
  {(1,e^{2\pi(s+it)},x)\sim
   (-1,e^{-2\pi(s+it)},(\phi_t^0)^{-1}(x))}.
\]
下文简记
\[
  z=e^{2\pi(s+it)},
\]
所以 $2\pi t$ 是 $z$ 的幅角。注意，底球面
$S^2=\mathbb C\cup\{\infty\}$ 是两份 $\mathbb C_\pm$ 之并，其中
$\mathbb C_+\equiv\mathbb C$，而
$\mathbb C_-\equiv(\mathbb C\setminus\{0\})\cup\{\infty\}$。
$\widetilde M^0$ 的一个截面是一对光滑映射 $u_\pm^0:\mathbb C\to M$，满足
\[
  u_+^0(z)=\phi_t^0\bigl(u_-^0(z^{-1})\bigr),
\]
$\widetilde M^\infty$ 的情形类似。为形成\BookRecordIndex{T0278}纤维和，把 $\widetilde M^0$ 在
$0\in\mathbb C_+$（原点）上的\BookRecordIndex{T0278}纤维和 $\widetilde M^\infty$ 在
$0\in\mathbb C_-$（无穷远点）上的纤维都与 $M$ 等同。然后从
$\widetilde M^0$、$\widetilde M^\infty$ 中分别删去这些纤维的邻域
$\mathbb C_+\times M$、$\mathbb C_-\times M$，再用复合映射
$\phi_t^\infty\circ\phi_t^0$，把 $\widetilde M^0$ 中余下的
$\mathbb C_-\times M$ 与 $\widetilde M^\infty$ 中余下的
$\mathbb C_+\times M$ 粘在一起，形成\BookRecordIndex{T0278}纤维和。这正是 $\widetilde M^{\infty0}$；它的一个
截面是一对光滑映射 $u_\pm:\mathbb C\to M$，满足
\[
  u_+(z)=\phi_t^\infty\circ\phi_t^0\bigl(u_-(z^{-1})\bigr).
\]

现在采取如下策略。从 $\widetilde M^0$ 和 $\widetilde M^\infty$ 中在公共纤维上相合的两个
$J$-全纯截面 $\widetilde u^0,\widetilde u^\infty$ 出发，构造
$\widetilde M^{\infty0}$ 的一族近似 $J^R$-全纯截面 $\widetilde u^R$（\BookRecordIndex{T0394}预粘合构造）。随后，
第 10 章的\BookRecordIndex{T0254}粘合论证将给出附近一族真正的 $J^R$-全纯截面。建立这一点后，像第~\ref{sec:10-8}
节那样论证即可完成证明。具体地，在 $S^2\setminus\{0\}$ 上固定 $\ell$ 个\BookRecordIndex{T0066}标记点
（用于 $\widetilde M^0$），在 $S^2\setminus\{\infty\}$ 上固定 $k-\ell$ 个\BookRecordIndex{T0066}标记点
（用于 $\widetilde M^\infty$），并在这 $k$ 个纤维中选取表示给定上同调类的适当子流形。
当 $R$ 充分大时，必须比较两种数目：其一是 $\widetilde M^{\infty0}$ 中穿过这 $k$ 个纤维
内给定子流形的 $J^R$-全纯截面数；其二是由 $\widetilde M^0$ 与
$\widetilde M^\infty$ 中在\BookRecordIndex{T0066}标记点处满足相同约束的 $J$-全纯截面组成的\BookRecordIndex{T0226}连通对的数目。像
第~\ref{sec:10-8} 节一样，这立即由\BookRecordIndex{T0256}粘合定理得到。反过来，\BookRecordIndex{T0256}粘合定理也由\BookRecordIndex{T0394}预粘合构造和
第 10 章第~\ref{sec:10-3}--\ref{sec:10-7} 节中完全相同的论证得到。

下面给出当前情形下\BookRecordIndex{T0394}预粘合构造的细节。考虑 $\widetilde M^0$ 的截面
$\widetilde u^0=(u_+^0,u_-^0)$ 与 $\widetilde M^\infty$ 的截面
$\widetilde u^\infty=(u_+^\infty,u_-^\infty)$，并假设
\[
  u_+^0(0)=u_-^\infty(0)=:x.
\]
对每个充分大的实数 $R$，定义 $\widetilde M^{\infty0}$ 的截面
$\widetilde u^R=(u_+^R,u_-^R)$，使得
\[
u_+^R(z)=
\begin{cases}
  u_+^\infty(R^2z),& |z|\le\dfrac{\delta}{2R},\\[3pt]
  \phi_t^\infty(x),& \dfrac{\delta}{R}\le|z|\le\dfrac1{\delta R},\\[3pt]
  \phi_t^\infty(u_+^0(z)),& |z|\ge\dfrac2{\delta R},
\end{cases}
\]
以及
\[
u_-^R(1/z)=
\begin{cases}
  (\phi_t^0)^{-1}\bigl(u_-^\infty(1/(R^2z))\bigr),
    & |z|\le\dfrac{\delta}{2R},\\[3pt]
  (\phi_t^0)^{-1}(x),
    & \dfrac{\delta}{R}\le|z|\le\dfrac1{\delta R},\\[3pt]
  u_-^0(1/z),& |z|\ge\dfrac2{\delta R}.
\end{cases}
\]
读者可以核对，只要 $u_+^R(z)$ 和 $u_-^R(1/z)$ 已定义，就对所有 $R$ 以及所有
$z=e^{2\pi(s+it)}$ 有
\[
  \phi_t^\infty\circ\phi_t^0\bigl(u_-^R(1/z)\bigr)=u_+^R(z).
\]
当 $R$ 充分大时，可以像第~\ref{sec:10-2} 节那样，通过插值把 $u_+^R$ 延拓到整个
$\mathbb C$：在\BookRecordIndex{T0170}环域 $\delta/(2R)\le|z|\le\delta/R$ 上，利用 $u_-^\infty$ 在
$u_+^\infty$ 与 $\phi_t^\infty(x)$ 之间插值；在
$1/(\delta R)\le|z|\le2/(\delta R)$ 上，利用 $u_+^0$ 在
$\phi_t^\infty(x)$ 与 $\phi_t^\infty(u_+^0)$ 之间插值。更准确地说，因为
$u_+^0(0)=u_-^\infty(0)=x$，存在光滑映射
$\zeta_+^0,\zeta_-^\infty:\mathbb C\to T_xM$，使得当 $|z|$ 充分小时
\[
  u_+^0(z)=\exp_x(\zeta_+^0(z)),
  \qquad
  u_-^\infty(z)=\exp_x(\zeta_-^\infty(z)),
  \qquad
  \zeta_+^0(0)=\zeta_-^\infty(0)=0.
\]
取光滑\BookRecordIndex{T0188}截断函数 $\rho:\mathbb C\to[0,1]$，使其在 $|z|\le1$ 时为零，在
$|z|\ge2$ 时为一，并定义
\[
\begin{aligned}
u_+^R(z):={}&\phi_t^\infty\!\left(
  \exp_x\!\left(
  \rho\!\left(\frac\delta{Rz}\right)
  \zeta_-^\infty\!\left(\frac1{R^2z}\right)
  +\rho(\delta Rz)\zeta_+^0(z)
  \right)\right),\\
u_-^R(1/z):={}&(\phi_t^0)^{-1}\!\left(
  \exp_x\!\left(
  \rho\!\left(\frac\delta{Rz}\right)
  \zeta_-^\infty\!\left(\frac1{R^2z}\right)
  +\rho(\delta Rz)\zeta_+^0(z)
  \right)\right)
\end{aligned}
\]
其中 $\delta/(2R)\le|z|\le R/(2\delta)$。第~\ref{sec:10-2} 节的
式~\eqref{eq:10-2-2}，对应于对所有 $t$ 均有
$\phi_t^0=\phi_t^\infty=\operatorname{id}$ 的特例。

这个\BookRecordIndex{T0394}预粘合构造定义一族映射 $f^R$，其定义域是由
$\widetilde M^0,\widetilde M^\infty$ 中分别表示类
$\widetilde A^0,\widetilde A^\infty$、并满足
$u_+^0(0)=u_-^\infty(0)$ 的 $J$-全纯截面对
$(\widetilde u^0,\widetilde u^\infty)$ 组成的\BookRecordIndex{T0246}模空间
\[
  \mathcal M(\widetilde A^0,\widetilde A^\infty;J^0,J^\infty).
\]
它把该\BookRecordIndex{T0246}模空间映到 $\widetilde M^{\infty0}$ 中表示同调类
\[
  \widetilde A^{\infty0}:=\widetilde A^\infty\#\widetilde A^0
\]
的截面空间。为了使\BookRecordIndex{T0254}粘合论证成立，必须像第 10 章那样假设\BookRecordIndex{T0246}模空间
$\mathcal M(\widetilde A^0,J^0)$、$\mathcal M(\widetilde A^\infty,J^\infty)$ 都正则，
并且\BookRecordIndex{T0283}求值映射
\[
\begin{aligned}
  \mathcal M(\widetilde A^0,J^0)\times
  \mathcal M(\widetilde A^\infty,J^\infty)
  &\longrightarrow M\times M,\\
  (\widetilde u^0,\widetilde u^\infty)
  &\longmapsto(u_+^0(0),u_-^\infty(0))
\end{aligned}
\]
横截于对角线。从这里开始，证明逐字等同于第 10 章的证明，故略去。这就证明了
定理~\ref{thm:11-4-6}。
\end{proof}

\begin{remark}\label{rem:11-4-7}
\begin{enumerate}[label=(\roman*)]
\item 如果把 \BookRecordIndex{T0046}Seidel 元 $\mathcal S(\phi)$ 定义为弗洛尔上同调上的算子，那么
定理~\ref{thm:11-4-3} 是一个同义反复。不过，这一叙述中非平凡的部分，是把该观察翻译回
当前语境，即证明弗洛尔上同调上的算子与上面定义的 \BookRecordIndex{T0047}Seidel 表示一致。这种对应的证明涉及
与定理~\ref{thm:11-4-3} 上述证明同类的\BookRecordIndex{T0256}粘合定理。进一步讨论见第 12.5 节。

\item 我们在上面根据分裂规则定理~\ref{thm:11-4-6} 作计算，证明了
定理~\ref{thm:11-4-3}。如果像注~\ref{rem:11-4-4} 那样，把 \BookRecordIndex{T0046}Seidel 元
$\mathcal S(\phi)$ 看成单位元 $1$ 在自同构 $\mathcal S(\phi)$ 下的像，就能给出该论证的
如下几何解释。与往常一样，用同调语言来说明最为容易。为此，记
$\mathcal S_*(\phi)$ 为 $\operatorname{QH}_*(M;\Lambda)$ 的自同构，它把 $w_1$ 上纤维
中的同调类 $\alpha$ 送到 $w_2$ 上纤维中的类
\[
  \mathcal S_*(\phi)(\alpha):=
  \sum_{\widetilde A\in H_2^\sigma(\widetilde M_\phi)}
  \sum_{i,j}
  \operatorname{GW}^{\widetilde M_\phi,\mathbf w}_{\widetilde A,2}
  (\alpha,\varepsilon_i)g^{ij}\varepsilon_j,
\]
其中 $\varepsilon_i:=\operatorname{PD}(e_i)$。换言之，我们用式~\eqref{eq:11-4-4}
右端的同调版本定义 $\mathcal S_*(\phi)(\alpha)$。于是
\[
  \mathcal S_*(\phi)(\alpha)
  =\operatorname{PD}\circ\mathcal S(\phi)\circ\operatorname{PD}(\alpha),
\]
其中
$\operatorname{PD}:\operatorname{QH}^*(M;\Lambda)\to
\operatorname{QH}_*(M;\check\Lambda)$ 是注~\ref{rem:11-1-20} 中定义的 Poincar\'e 对偶
同构。我们的目标是证明复合规则
\[
  \mathcal S_*(\phi^\infty\phi^0)
  =\mathcal S_*(\phi^\infty)\circ\mathcal S_*(\phi^0).
\]
对 $a=0,\infty$，把 $\mathcal S_*(\phi^a)$ 看成一个自同构：它把
$\widetilde M^a$ 在 $0\in\mathbb C_-$（无穷远点）上的纤维 $M_\infty^a$ 的\BookRecordIndex{T0230}量子同调，
送到其在 $0\in\mathbb C_+$（原点）上的纤维 $M_0^a$ 的\BookRecordIndex{T0230}量子同调。循环
$\alpha\subset M_\infty^0$ 的像 $\mathcal S_*(\phi^0)(\alpha)$，由
$\widetilde M^0$ 中各截面与 $M_0^0$ 的交表示；所以复合
$\mathcal S_*(\phi^\infty)\circ\mathcal S_*(\phi^0)(\alpha)$ 的像，由在公共纤维上相合的
截面对与 $M_0^\infty$ 的交表示。另一方面，$\mathcal S_*(\phi^\infty\phi^0)(\alpha)$
按定义通过计数 $\widetilde M^{\infty0}$ 的截面得到。因此，计算化为
定理~\ref{thm:11-4-6} 的\BookRecordIndex{T0254}粘合陈述。在这一框架中还必须证明 $\mathcal S_*(\phi)$ 是模同态，
也就是
\[
  \mathcal S_*(\phi)(\alpha)
  =\mathcal S_*(\phi)(\alpha*1)
  =\alpha*\mathcal S_*(\phi)(1)
  =\mathcal S_*(\phi)*\alpha.
\]
这可仿照定理~\ref{thm:11-4-3} 证明的第一步，并取其中一个环路为平凡环路来得到。
文献 \srcite{228,269} 采取的正是这一观点。
\end{enumerate}
\end{remark}

\begin{example}\label{ex:11-4-8}
在多数已经算出 \BookRecordIndex{T0046}Seidel 元 $\mathcal S(\psi)\in\operatorname{QH}^*(M)$ 的情形中，环路
$\psi$ 由 $(M,\omega)$ 上的一个圆周作用表示。令 $K:M\to\mathbb R$ 是其\BookRecordIndex{T0204}矩映射，并作
归一化使
\[
  \int_M K\omega^n=0,
\]
再把相应环路记作 $\lambda_K$。则 $K$ 取得最大值的点组成一个连通辛子流形
$F_K\subset M$，且该子流形由作用固定：
\[
  K(F_K)=\max_{x\in M}K(x).
\]
进一步假设这个最大值分支是简单的，即圆周在 $F_K$ 的一个去心邻域
$U\setminus F_K$ 上自由作用。McDuff--Tolman \srcite{279} 证明了
\begin{equation}\label{eq:11-4-7}
  \mathcal S(\lambda_K)
  =\operatorname{PD}(F_K)\otimes
   q^{-\operatorname{codim}(F_K)/2}t^{\sup K}
  +\sum_{\widetilde\omega(\widetilde A)>\sup K}
   a_{\widetilde A}\otimes
   q^{c(\widetilde A)}t^{\widetilde\omega(\widetilde A)}.
\end{equation}
因此，$\mathcal S(\lambda_K)$ 的首项由几何以很简单的方式确定。举例来说，考虑由高度函数
$K$ 生成的 $S^2$ 旋转。这时
\[
  \mathcal S(\lambda_K)=\operatorname{PD}([\mathrm{pt}])\otimes q^{-1}t^\mu,
  \qquad
  \mu=\sup K=\frac12\int_{S^2}\omega.
\]
在若干有趣情形中，式~\eqref{eq:11-4-7} 有非零低阶项（参见 Seidel \srcite{363}）。
不过，文献 \srcite{279} 证明：若 $(M,\omega,\lambda_K)$ 是 Fano 的，也就是说，它有一个
不允许非正陈数 $J$-全纯球面的不变 $\omega$-\BookRecordIndex{T0372}驯服\BookRecordIndex{T0088}殆复结构；并且\BookRecordIndex{T0418}最小陈数 $N$ 至少是
$F_K$ 余维的一半，则这些项不会出现。例如，若 $\lambda_K$ 是 Fano 环面簇上固定一个面片的
圆周作用之一，这些项就不会出现。
\end{example}

\begin{example}\label{ex:11-4-9}
回到命题~\ref{prop:9-7-2} 中考虑的圆周作用。它按公式
\[
  \psi_t(z,w)=(z,\Phi_{z,t}(w))
\]
作用在 $S^2\times S^2$ 上，因而旋转到第一个球面的投影之纤维，并固定对角线与反对角线。
第 9 章给出了两个证明，说明当 $\lambda>1$ 时，$S^2\times S^2$ 上存在一个
$S^1$-不变\BookRecordIndex{T0363}辛形式，它\BookRecordIndex{T0361}辛同痕于标准形式
$\omega_\lambda:=\lambda\pi_1^*\sigma+\pi_2^*\sigma$。还可改写
例~\ref{ex:11-3-12} 的构造，给出
包含这一 $S^1$ 作用的 $S^2\times S^2$ 环面结构，从而得到第三个证明：只须把
$w_2$ 取作向量 $(2,1)$ 而非 $(1,1)$，并适当选择 $\tau$。命题~\ref{prop:9-7-2}
中更有趣的部分断言：当 $\lambda\le1$ 时，不可能把哈密顿环路 $\{\psi_t\}$ 同伦到
$(S^2\times S^2,\omega_\lambda)$ 的\BookRecordIndex{T0362}辛微分同胚群中。现在概述另一个证明。由连续性，只须
考虑 $\lambda=1$。

我们在例~\ref{ex:11-1-13} 中看到，单调情形的\BookRecordIndex{T0229}量子上同调环
$\operatorname{QH}^*(S^2\times S^2;\Lambda)$ 可构造成有零因子的环，而当 $\lambda>1$ 时
却没有零因子。取当前的 Novikov 环
$\Lambda_\lambda:=\Lambda^{\mathrm{univ}}[q,q^{-1}]$，这一现象非常清楚。记
\[
  A:=[S^2\times\mathrm{pt}],\qquad
  B:=[\mathrm{pt}\times S^2],\qquad
  a:=[\pi_1^*\sigma]=\operatorname{PD}(B),\qquad
  b:=[\pi_2^*\sigma]=\operatorname{PD}(A).
\]
则同态 $\varphi:\Gamma(M,\omega)\to\Lambda$ 给为
\[
  \varphi(A):=qt^\lambda,
  \qquad
  \varphi(B):=qt.
\]
因此，当 $\lambda=1$ 时，$a*a=qt$ 与 $b*b=qt^\lambda$ 相等。进一步，
\[
  (a+b)(a-b)=qt(1-t^{\lambda-1}).
\]
所以当 $\lambda>1$ 时，$a+b$ 是
$\operatorname{QH}^*(S^2\times S^2;\Lambda_\lambda)$ 中的单位，而当 $\lambda=1$ 时，
它是零因子。

圆周作用 $\{\psi_t\}$ 的最大不动点分支 $F_K$ 是对角线，其 Poincar\'e 对偶是
$H^2(S^2\times S^2)$ 中的类 $a+b$。Seidel \srcite{363} 通过直接计算证明
$\mathcal S(\psi)$ 中没有低阶项；这也由例~\ref{ex:11-4-8} 的讨论得到。因此，当
$\lambda>1$ 时，$\mathcal S(\psi)$ 形如 $(a+b)\otimes\mu_\lambda$，其中
$\mu_\lambda\in\Lambda_\lambda$。\BookRecordIndex{T0046}Seidel 元由 $\widetilde M_\psi$ 中的\BookRecordIndex{T0143}格罗莫夫--威滕
不变量给出，而这些不变量不依赖于 $M$ 上\BookRecordIndex{T0363}辛形式的形变类。因此，如果 $\psi$ 同伦于
$\operatorname{Symp}(M,\omega_1)$ 中的一个环路，其 \BookRecordIndex{T0046}Seidel 元
$\mathcal S(\psi)\in\operatorname{QH}^*(S^2\times S^2;\Lambda_1)$ 仍会形如
$(a+b)\otimes\mu_1$，其中 $\mu_1\in\Lambda_1$。由于这不是单位，该同伦不可能存在。
\end{example}

\paragraph{光滑环面簇的小\BookRecordIndex{T0229}量子上同调。}
现在回到第~\ref{subsec:11-3-1} 小节开始的\BookRecordIndex{T0169}环面流形小\BookRecordIndex{T0229}量子上同调环
$\operatorname{QH}^*(M)$ 的讨论。我们的叙述紧随 McDuff--Tolman \srcite{279}。假设
$(M,\omega)$ 是闭辛 $2n$-流形，并带有 $n$-环面 $T$ 的有效哈密顿群作用。令
$\mu:M\to\mathfrak t^*$ 是这一作用的\BookRecordIndex{T0204}矩映射，并记其像为 $\Delta:=\mu(M)$。回忆
$\Delta$ 是 $\mathfrak t^*$ 中的凸多面体，可写成
\[
  \Delta=\left\{\eta\in\mathfrak t^*\ \middle|\
  \langle\eta,\bar e_\nu\rangle\le\kappa_\nu\right\},
\]
其中 $\bar e_1,\ldots,\bar e_N\in\mathfrak t$ 属于整数格，且
$\kappa_1,\ldots,\kappa_N\in\mathbb R$（这里从面片的内法向量改用外法向量）。再回忆诸面
\[
  \Delta_I:=\left\{\eta\in\Delta\ \middle|\
  \nu\in I\Longrightarrow\langle\eta,\bar e_\nu\rangle=\kappa_\nu\right\},
  \qquad I\subset I_0:=\{1,\ldots,N\}.
\]
由式~\eqref{eq:11-3-15}，对每个\BookRecordIndex{T0060}本原子集 $I\subset I_0$（即
$\Delta_I=\varnothing$，而且只要 $I'\subsetneq I$ 就有
$\Delta_{I'}\ne\varnothing$），都存在唯一向量
$d=(d_1,\ldots,d_N)\in\mathbb Z^N$，满足
\begin{equation}\label{eq:11-4-8}
  \sum_{\nu=1}^N d_\nu\bar e_\nu=0,
  \qquad
  d_\nu=1\quad(\nu\in I),
  \qquad
  d_\nu\le0\quad(\nu\notin I).
\end{equation}
每个这样的整数向量 $d$ 对应一个同调类 $A_d\in H_2(M)$，其第一陈数和辛面积分别为
\begin{equation}\label{eq:11-4-9}
  c_1(A_d)=\sum_{\nu=1}^N d_\nu,
  \qquad
  \omega(A_d)=\sum_{\nu=1}^N\kappa_\nu d_\nu.
\end{equation}
（见式~\eqref{eq:11-3-18}，其中把 $\kappa_\nu$ 换成 $-\kappa_\nu$。）我们希望理解\BookRecordIndex{T0229}量子
上同调中与通常上同调关系 $\prod_{\nu\in I}u_\nu=0$ 对应的关系，其中 $I$ 是任意本原
指标集。

下面直接从几何上解释 Stanley--Reisner 关系的量子形变。因为面 $\Delta_\nu$ 的对偶向量
$\bar e_\nu$ 是整数向量，它在 $M$ 上诱导圆周作用 $\lambda_\nu$，即 $T$ 中与
$\bar e_\nu$ 相切的圆所给出的作用。它的\BookRecordIndex{T0204}矩映射是
$\langle\mu,\bar e_\nu\rangle:M\to\mathbb R$。所以关系
$\sum_{\nu=1}^N d_\nu\bar e_\nu=0$ 断言，$\pi_1(\operatorname{Ham}(M,\omega))$
中相应的乘积是常值环路：
\[
  \prod_{\nu=1}^N(\lambda_\nu)^{d_\nu}=1.
\]
由定理~\ref{thm:11-4-3}，该乘积在 \BookRecordIndex{T0047}Seidel 表示下的像必须是单位元。换言之，
\begin{equation}\label{eq:11-4-10}
  \prod_{\nu\in I}\mathcal S(\lambda_\nu)
  =\prod_{\nu\notin I}\mathcal S(\lambda_\nu)^{|d_\nu|}
  \in\operatorname{QH}^*(M;\Lambda).
\end{equation}
这里的乘积均在\BookRecordIndex{T0229}量子上同调中理解。如例~\ref{ex:11-4-8} 所指出的，在 Fano 情形中，
$\mathcal S(\lambda_\nu)$ 的表达式没有低阶项。此外，在所有情形中，第 $\nu$ 个圆周作用的
\BookRecordIndex{T0204}矩映射取得最大值的子流形 $X_\nu\subset M$，恰是第 $\nu$ 个面片在\BookRecordIndex{T0204}矩映射下的原像，故其
Poincar\'e 对偶为 $\bar w_\nu\in H^2(M)$：
\[
  \kappa_\nu=\sup_{x\in M}\langle\mu,\bar e_\nu\rangle
  =\langle\mu(X_\nu),\bar e_\nu\rangle,
  \qquad
  X_\nu:=\mu^{-1}(\Delta_\nu).
\]
所以例~\ref{ex:11-4-8} 给出
\[
  \mathcal S(\lambda_\nu)=\bar w_\nu\otimes q^{-1}t^{\kappa_\nu},
  \qquad
  \bar w_\nu:=\operatorname{PD}([X_\nu])\in H^2(M).
\]
这证明了定理~\ref{thm:11-3-4}（取式~\eqref{eq:11-3-17} 的版本），其中\BookRecordIndex{T0231}量子系数环为
\[
  \Lambda:=\Lambda^{\mathrm{univ}}[q,q^{-1}].
\]
这还不是完整结果，因为它只保留了同调类 $A\in H_2(M)$ 的陈数和面积信息。不过，下一则
注说明了如何证明完整结果。

\begin{remark}\label{rem:11-4-10}
因为 $M$ 是 Fano 的，可取多项式环
$\Lambda_0:=\mathbb Q[q_1,\ldots,q_k]$ 作为其\BookRecordIndex{T0230}量子同调的系数。对 $M$ 上每个圆周作用
$\lambda_\nu$，$\widetilde M_{\lambda_\nu}$ 中都有典范截面 $\sigma_\nu$，由点
$\{\pm1\}\times\mathbb C\times\{x\}$ 组成，其中 $x\in X_\nu$。类似地，对任意
$d\in\mathbb Z$，给定圆周作用 $\lambda_\nu^d$，把 $\widetilde M_{\lambda_\nu^d}$ 中\BookRecordIndex{T0204}矩映射
$dK_\nu$ 取得最大值的截面定义为 $(\sigma_\nu)^d$。文献 \srcite{279} 中的计算表明，
定义~\ref{def:11-4-1} 的 \BookRecordIndex{T0046}Seidel 元
$\mathcal S(\lambda_\nu,\sigma_\nu)\in\operatorname{QH}^*(M;\Lambda_0)$ 等于
$\bar w_\nu$。考虑上述关系 $\sum d_\nu\bar e_\nu=0$。由
定理~\ref{thm:11-4-6}，
\[
  \prod_{\nu=1}^N
  \mathcal S\bigl(\lambda_\nu^{d_\nu},(\sigma_\nu)^{d_\nu}\bigr)
  =\mathcal S(0,\sigma_d)=1\otimes e^A,
\]
其中 $\sigma_d$ 是把诸截面 $(\sigma_\nu)^{d_\nu}$ 串接而成的平凡丛截面，且属于
$H_2(S^2\times M)$ 中的类 $[S^2\times\mathrm{pt}]+A$。这说明
$\operatorname{QH}^*(M;\Lambda_0)$ 中有形如
\[
  \prod_{\nu\in I}\bar w_\nu
  =e^A\prod_{\nu\notin I}(\bar w_\nu)^{|d_\nu|}
\]
的关系。还须把 $A$ 与之前记作 $A_d$ 的元素等同。但对 $[\omega]$ 的每个选择，这两个类
在 $\Lambda^{\mathrm{univ}}[q,q^{-1}]$ 中的像都相同。由于 $[\omega]$ 可在 $H^2(M)$ 的一个
开子集中变化，遂有 $A=A_d$。

这一方法的优点在于，不论\BookRecordIndex{T0169}环面流形 $(M,\omega)$ 是否 Fano，关系~\eqref{eq:11-4-10}
总是成立。事实上，由例~\ref{ex:11-4-8}，\BookRecordIndex{T0046}Seidel 元总可写成
\[
  \mathcal S(\lambda_\nu)=y_\nu\otimes q^{-1}t^{\kappa_\nu},
  \qquad
  y_\nu=\bar w_\nu+\sum_k\sum_{\varepsilon>0}a_{k,\varepsilon}q^kt^\varepsilon.
\]
\sourcefootnote{原书把 $t$ 的指数印作本节从未定义的 $c_\nu$；紧前面的\BookRecordIndex{T0204}矩映射最大值公式及
\BookRecordIndex{T0046}Seidel 元首项均唯一给出该指数为 $\kappa_\nu$。}
如第~\ref{subsec:11-3-1} 小节所述，有满环同态
\[
  \frac{\Lambda[u_1,\ldots,u_N]}{P(\Delta)}
  \longrightarrow\operatorname{QH}^*(M;\Lambda),
  \qquad
  u_\nu\longmapsto\bar w_\nu,
  \qquad
  \Lambda:=\Lambda^{\mathrm{univ}}[q,q^{-1}].
\]
所以 $y_\nu\in\operatorname{QH}^*(M;\Lambda)$ 可提升为某个形如
$Y_\nu=u_\nu+\text{高阶项}$ 的元素
$Y_\nu\in\Lambda[u_1,\ldots,u_N]$。利用这些多项式，可以定义量子
Stanley--Reisner 理想 $SR_Y(\Delta)\subset\Lambda[u_1,\ldots,u_N]$：
\begin{equation}\label{eq:11-4-11}
  SR_Y(\Delta):=
  \left\langle
  \prod_{\nu\in I}Y_\nu
  =q^{c_1(A_d)}t^{\omega(A_d)}
   \prod_{\nu\notin I}Y_\nu^{|d_\nu|}
  \ \middle|\ I\text{ 是本原的}
  \right\rangle.
\end{equation}
这里 $d_\nu=d_\nu(I)$ 如式~\eqref{eq:11-4-8}，而 $c_1(A_d)$、$\omega(A_d)$ 由
式~\eqref{eq:11-4-9} 给出。注意，$SR_Y(\Delta)$ 依赖于 $Y_\nu$；一般而言，不预先知道
$\operatorname{QH}^*(M;\Lambda)$ 的环结构，就无法描述这些 $Y_\nu$。所以 $M$ 的小\BookRecordIndex{T0229}量子
上同调形如
\[
  \operatorname{QH}^*(M;\Lambda)
  \cong\frac{\Lambda[u_1,\ldots,u_N]}{P(\Delta)+SR_Y(\Delta)},
  \qquad
  \Lambda:=\Lambda^{\mathrm{univ}}[q,q^{-1}].
\]
\end{remark}

\begin{example}[NEF 环面簇]\label{ex:11-4-11}
假设 $M$ 是 NEF 的，即每个有全纯代表元的类 $B\in H_2(M)$ 都满足 $c_1(B)\ge0$。
在此情形中，\BookRecordIndex{T0046}Seidel 元 $y_\nu$ 里可能有高阶项。不过，文献 \srcite{279} 证明，这时
$\mathcal S(\lambda_\nu)$ 公式中出现的上同调类 $a_{k,\varepsilon}$ 次数都是零或二。因此，
它们自然提升到
$\Lambda[u_1,\ldots,u_N]$\sourcefootnote{原书把此处多项式环的末指标印作 $n$；但
本段由矩多面体的 $N$ 个面片定义变量 $u_1,\ldots,u_N$，前文环及
式~\eqref{eq:11-4-11} 也都以 $u_N$ 为末项，故补正为 $N$。}，而无须预先知道量子乘法。
所以，一旦算出 \BookRecordIndex{T0046}Seidel 元 $\mathcal S(\lambda_\nu)$，便可写出式~\eqref{eq:11-4-11}
中的\BookRecordIndex{T0232}量子 Stanley--Reisner 环。这一结果，尤其是把 Stanley--Reisner 环 $SR_Y$ 中的
$u_\nu$ 换成 $Y_\nu$，与 \srcite{76, 11.2.5.2} 所讨论的 Givental 变量替换公式有关。
Givental 镜像定理提供了计算 \BookRecordIndex{T0046}Seidel 元 $Y_i$ 的方法：它联系函数 $\widetilde I$ 与
$\widetilde J$，前者可由组合数据计算，后者则决定\BookRecordIndex{T0229}量子上同调中的关系；见
Gonzalez--Iritani \srcite{155}。
\end{example}
